% GENERATED FILE. Edit canonical lessons, then run npm run generate:books.
% canonical-source-hash: fnv1a64:7fac7ed7d2054b0f
% canonical-lessons: ML-C400-satru, ML-C400-yoddhavu, ML-C400-caran, ML-C400-tolppikkuka, ML-C400-kilatannuka

\chapter{Enemies and Spies}
\label{ch:ml-enemies-and-spies}
\hlchaptermodality{400}

\begin{chapteropening}
\textbf{By the end of this chapter:} \emph{I can say an enemy, a warrior, a spy, to beat in battle and to surrender.}

Complete the last lesson of chapter 400: I can say an enemy, a warrior, a spy, to beat in battle and to surrender.
\end{chapteropening}

\section[\texorpdfstring{\ml{ശത്രു} (śatru)}{ശത്രു (śatru)}]{\ml{ശത്രു} (śatru) — an enemy}

\label{lesson:ML-C400-satru}



\begin{quote}
Before the new one: say the Malayalam for a soldier, then the Malayalam for to conquer.
\end{quote}

\subsection*{You'll want to know: \ml{ശത്രു}}
\textbf{\ml{ശത്രു}} — \emph{śatru} — \textquotedblleft{}an enemy\textquotedblright{}.

\begin{quote}
\ml{ശത്രുവിന്റെ} \ml{സൈന്യം} \ml{കോട്ട} \ml{വളഞ്ഞു}.
\end{quote}

\begin{quote}
\emph{śatruvinṟe sainyaṁ kōṭṭa vaḷaññu.}
\end{quote}

\begin{quote}
The enemy's army surrounded the fort.
\end{quote}

Plural: \ml{ശത്രുക്കൾ}, \emph{śatrukkaḷ}.

Plural \ml{ശത്രുക്കൾ}.

\begin{cousinweb}
\textbf{\ml{ശത്രു}} (\emph{śatru}) is borrowed from Sanskrit \emph{śatru}, \textquotedblleft{}an enemy\textquotedblright{} (CDIAL 12286).
\end{cousinweb}

\begin{grammarlens}[title={What you've built}]
One more word for telling what happened.
\end{grammarlens}

\subsection*{Your turn}
\begin{itemize}
\raggedright
  \item \emph{Say it:} \emph{śatru}
  \item \emph{Say it:} \emph{śatru}, once more
  \item \emph{Say it:} \emph{śatruvinṟe sainyaṁ kōṭṭa vaḷaññu.} — The enemy's army surrounded the fort.
  \item \emph{Recall:} say the Malayalam for a soldier, then the Malayalam for to conquer, then say \emph{śatru} again
\end{itemize}

\begin{tcolorbox}[breakable,colback=teal!4,colframe=teal!35!black,title={Before you move on}]
What does \ml{ശത്രു} mean? (\textquotedblleft{}An enemy\textquotedblright{}.) Say it once more.
\end{tcolorbox}
\section[\texorpdfstring{\ml{യോദ്ധാവ്} (yōddhāvŭ)}{യോദ്ധാവ് (yōddhāvŭ)}]{\ml{യോദ്ധാവ്} (yōddhāvŭ) — a warrior}

\label{lesson:ML-C400-yoddhavu}



\begin{quote}
Before the new one: say the Malayalam for to conquer, then the Malayalam for an enemy.
\end{quote}

\subsection*{You'll want to know: \ml{യോദ്ധാവ്}}
\textbf{\ml{യോദ്ധാവ്}} — \emph{yōddhāvŭ} — \textquotedblleft{}a warrior\textquotedblright{}.

\begin{quote}
\ml{ആ} \ml{യോദ്ധാവ്} \ml{ഒറ്റയ്ക്ക്} \ml{ശത്രുക്കളോട്} \ml{യുദ്ധം} \ml{ചെയ്തു}.
\end{quote}

\begin{quote}
\emph{ā yōddhāvŭ oṟṟaykkŭ śatrukkaḷōṭŭ yuddhaṁ ceytu.}
\end{quote}

\begin{quote}
That warrior fought the enemies alone.
\end{quote}

From the same root as \ml{യുദ്ധം}, war.

\begin{cousinweb}
\textbf{\ml{യോദ്ധാവ്}} (\emph{yōddhāvŭ}) is borrowed from Sanskrit \emph{yoddhṛ}, \textquotedblleft{}a fighter, a warrior, a soldier\textquotedblright{} (Monier-Williams), from the root \emph{yudh}, \textquotedblleft{}to fight\textquotedblright{}, the root of \ml{യുദ്ധം} (\emph{yuddhaṁ}).
\end{cousinweb}

\begin{grammarlens}[title={What you've built}]
One more word for telling what happened.
\end{grammarlens}

\subsection*{Your turn}
\begin{itemize}
\raggedright
  \item \emph{Say it:} \emph{yōddhāvŭ}
  \item \emph{Say it:} \emph{yōddhāvŭ}, once more
  \item \emph{Say it:} \emph{ā yōddhāvŭ oṟṟaykkŭ śatrukkaḷōṭŭ yuddhaṁ ceytu.} — That warrior fought the enemies alone.
  \item \emph{Recall:} say the Malayalam for to conquer, then the Malayalam for an enemy, then say \emph{yōddhāvŭ} again
\end{itemize}

\begin{tcolorbox}[breakable,colback=teal!4,colframe=teal!35!black,title={Before you move on}]
What does \ml{യോദ്ധാവ്} mean? (\textquotedblleft{}A warrior\textquotedblright{}.) Say it once more.
\end{tcolorbox}
\section[\texorpdfstring{\ml{ചാരൻ} (cāran)}{ചാരൻ (cāran)}]{\ml{ചാരൻ} (cāran) — a spy}

\label{lesson:ML-C400-caran}



\begin{quote}
Before the new one: say the Malayalam for an enemy, then the Malayalam for a warrior.
\end{quote}

\subsection*{You'll want to know: \ml{ചാരൻ}}
\textbf{\ml{ചാരൻ}} — \emph{cāran} — \textquotedblleft{}a spy\textquotedblright{}.

\begin{quote}
\ml{രാജാവ്} \ml{ശത്രുവിന്റെ} \ml{കൊട്ടാരത്തിലേക്ക്} \ml{ഒരു} \ml{ചാരനെ} \ml{അയച്ചു}.
\end{quote}

\begin{quote}
\emph{rājāvŭ śatruvinṟe koṭṭārattilēkkŭ oru cārane ayaccu.}
\end{quote}

\begin{quote}
The king sent a spy to the enemy's palace.
\end{quote}

One who goes about and watches.

\begin{cousinweb}
\textbf{\ml{ചാരൻ}} (\emph{cāran}) is borrowed from Sanskrit \emph{cāra}, \textquotedblleft{}one who moves about; a spy\textquotedblright{} (CDIAL 4754), from \emph{car}, \textquotedblleft{}to move\textquotedblright{}, with the Malayalam masculine ending -\ml{ൻ}.
\end{cousinweb}

\begin{grammarlens}[title={What you've built}]
One more word for telling what happened.
\end{grammarlens}

\subsection*{Your turn}
\begin{itemize}
\raggedright
  \item \emph{Say it:} \emph{cāran}
  \item \emph{Say it:} \emph{cāran}, once more
  \item \emph{Say it:} \emph{rājāvŭ śatruvinṟe koṭṭārattilēkkŭ oru cārane ayaccu.} — The king sent a spy to the enemy's palace.
  \item \emph{Recall:} say the Malayalam for an enemy, then the Malayalam for a warrior, then say \emph{cāran} again
\end{itemize}

\begin{tcolorbox}[breakable,colback=teal!4,colframe=teal!35!black,title={Before you move on}]
What does \ml{ചാരൻ} mean? (\textquotedblleft{}A spy\textquotedblright{}.) Say it once more.
\end{tcolorbox}
\section[\texorpdfstring{\ml{തോൽപ്പിക്കുക} (tōlppikkuka)}{തോൽപ്പിക്കുക (tōlppikkuka)}]{\ml{തോൽപ്പിക്കുക} (tōlppikkuka) — to beat in battle; to defeat}

\label{lesson:ML-C400-tolppikkuka}



\begin{quote}
Before the new one: say the Malayalam for a warrior, then the Malayalam for a spy.
\end{quote}

\subsection*{You'll want to know: \ml{തോൽപ്പിക്കുക}}
\textbf{\ml{തോൽപ്പിക്കുക}} — \emph{tōlppikkuka} — \textquotedblleft{}to beat in battle; to defeat\textquotedblright{}.

\begin{quote}
\ml{ചെറിയ} \ml{സൈന്യം} \ml{വലിയ} \ml{സൈന്യത്തെ} \ml{തോൽപ്പിച്ചു}.
\end{quote}

\begin{quote}
\emph{ceṟiya sainyaṁ valiya sainyatte tōlppiccu.}
\end{quote}

\begin{quote}
The small army defeated the big army.
\end{quote}

Past: \ml{തോൽപ്പിച്ചു}, \emph{tōlppiccu}.

The causative of \ml{തോൽക്കുക}, to lose: to make someone lose.

\begin{cousinweb}
\textbf{\ml{തോൽപ്പിക്കുക}} (\emph{tōlppikkuka}) is native Dravidian, the causative of \ml{തോൽക്കുക} (\emph{tōlkkuka}), \textquotedblleft{}to be defeated, to lose\textquotedblright{} (DEDR 3558).
\end{cousinweb}

\begin{grammarlens}[title={What you've built}]
One more word for telling what happened.
\end{grammarlens}

\subsection*{Your turn}
\begin{itemize}
\raggedright
  \item \emph{Say it:} \emph{tōlppikkuka}
  \item \emph{Say it:} \emph{tōlppikkuka}, once more
  \item \emph{Say it:} \emph{ceṟiya sainyaṁ valiya sainyatte tōlppiccu.} — The small army defeated the big army.
  \item \emph{Recall:} say the Malayalam for a warrior, then the Malayalam for a spy, then say \emph{tōlppikkuka} again
\end{itemize}

\begin{tcolorbox}[breakable,colback=teal!4,colframe=teal!35!black,title={Before you move on}]
What does \ml{തോൽപ്പിക്കുക} mean? (\textquotedblleft{}To beat in battle; to defeat\textquotedblright{}.) Say it once more.
\end{tcolorbox}
\section[\texorpdfstring{\ml{കീഴടങ്ങുക} (kīḻaṭaṅṅuka)}{കീഴടങ്ങുക (kīḻaṭaṅṅuka)}]{\ml{കീഴടങ്ങുക} (kīḻaṭaṅṅuka) — to surrender}

\label{lesson:ML-C400-kilatannuka}



\begin{quote}
Before the new one: say the Malayalam for a spy, then the Malayalam for to beat in battle.
\end{quote}

\subsection*{You'll want to know: \ml{കീഴടങ്ങുക}}
\textbf{\ml{കീഴടങ്ങുക}} — \emph{kīḻaṭaṅṅuka} — \textquotedblleft{}to surrender\textquotedblright{}.

\begin{quote}
\ml{ഒടുവിൽ} \ml{ശത്രുസൈന്യം} \ml{കീഴടങ്ങി}.
\end{quote}

\begin{quote}
\emph{oṭuvil śatrusainyaṁ kīḻaṭaṅṅi.}
\end{quote}

\begin{quote}
At last the enemy army surrendered.
\end{quote}

Past: \ml{കീഴടങ്ങി}, \emph{kīḻaṭaṅṅi}.

Also to give oneself up to the police: \ml{പ്രതി} \ml{പോലീസിൽ} \ml{കീഴടങ്ങി}.

\begin{cousinweb}
\textbf{\ml{കീഴടങ്ങുക}} (\emph{kīḻaṭaṅṅuka}) is a native Dravidian compound of \ml{കീഴ്}, \textquotedblleft{}below\textquotedblright{} (DEDR 1619), and \ml{അടങ്ങുക}, \textquotedblleft{}to submit, to yield\textquotedblright{} (DEDR 63): the intransitive partner of \ml{കീഴടക്കുക} (\emph{kīḻaṭakkuka}).
\end{cousinweb}

\begin{grammarlens}[title={What you've built}]
One more word for telling what happened.
\end{grammarlens}

\subsection*{Your turn}
\begin{itemize}
\raggedright
  \item \emph{Say it:} \emph{kīḻaṭaṅṅuka}
  \item \emph{Say it:} \emph{kīḻaṭaṅṅuka}, once more
  \item \emph{Say it:} \emph{oṭuvil śatrusainyaṁ kīḻaṭaṅṅi.} — At last the enemy army surrendered.
  \item \emph{Recall:} say the Malayalam for a spy, then the Malayalam for to beat in battle, then say \emph{kīḻaṭaṅṅuka} again
\end{itemize}

\begin{tcolorbox}[breakable,colback=teal!4,colframe=teal!35!black,title={Before you move on}]
What does \ml{കീഴടങ്ങുക} mean? (\textquotedblleft{}To surrender\textquotedblright{}.) Say it once more.
\end{tcolorbox}
